\chapter{Series and interchange}

\begin{orient}
\textbf{What this chapter is for.} When no antiderivative exists, expand. A geometric expansion of $1/(1\pm x)$, a binomial expansion, or a plain Maclaurin series converts one impossible integral into infinitely many possible ones, and the sum of those is the answer. This is why the closed forms in hard definite integration are so often $\zeta(2)$, $\zeta(3)$, $\ln2$, Catalan's constant $G$, or a polylogarithm: those are simply what $\sum n^{-s}$ and its relatives are called. The chapter gives the expand and integrate termwise protocol, the small dictionary of moments that lets you do the general term, the dictionary of constants that lets you recognise the sum, and, at equal length, the theorems that license the interchange. That last part is not ceremony. There is a worked example here where an unjustified swap produces a finite answer of the wrong sign, and another where two orders of a double integral give $\pi/4$ and $-\pi/4$.
\end{orient}

\section{Expand and integrate termwise}

The strategy is a trade. You give up a closed form for the integrand and accept an infinite sum, betting that each individual term is elementary and that the sum is a constant you can name. The bet pays off surprisingly often, because the integrands that resist antidifferentiation are exactly the ones built out of $\ln$, $\arctan$ and $1/(1\pm x^{n})$, and those have clean series with rational coefficients.

The protocol has five steps, and the fifth is the one that matters.

\begin{keynote}[The five steps]
\textbf{1.} Identify the factor to expand and the region where the expansion is valid.
\textbf{2.} Expand it, keeping the index bookkeeping explicit.
\textbf{3.} Integrate the general term in closed form. Almost always this is the log power moment $\int_0^1x^{a}\ln^{b}x\dd x$ or the Gamma moment $\int_0^\infty x^{a}\eu^{-nx}\dd x$.
\textbf{4.} Sum the resulting series and name the constant.
\textbf{5.} Justify the interchange $\int\sum=\sum\int$ by monotone convergence, dominated convergence, or Tonelli, and check the answer numerically.
\end{keynote}

There is a structural reason why the constants come out the way they do. Termwise integration on $[0,1]$ against a power $x^{n}$ produces $\frac{1}{n+1}$, and against $x^{n}\ln^{b}x$ produces $\frac{b!}{(n+1)^{b+1}}$ up to sign. So a single expansion delivers $\sum n^{-1}$ shapes, an expansion with one logarithm present delivers $\sum n^{-2}$, and two logarithms deliver $\sum n^{-3}$. Every logarithm in the integrand is worth one unit of weight in the answer. Counting logarithms before starting tells you which zeta value to expect, and an answer of the wrong weight is a sign of an arithmetic slip rather than a surprising identity.

The same accounting works on the half line, where the moment is $\int_0^\infty x^{s-1}\eu^{-nx}\dd x=\Gamma(s)n^{-s}$: the power of $x$ in the integrand sets the weight directly. This is why $\int_0^\infty\frac{x^{3}}{\eu^{x}-1}\dd x$ is a $\zeta(4)$ problem and could not have been anything else.

Geometric expansion is where to start, because it is the one expansion whose terms are all of one sign on the natural domain, which makes step 5 free.

\tech{Geometric expansion}{core}
\cue $\dfrac{1}{1\pm x^{n}}$ on $[0,1]$, or $\dfrac{1}{1-(xy)^{a}}$ on the unit square, or a Bose factor $\dfrac{1}{\eu^{x}-1}$, $\dfrac{1}{\eu^{x}+1}$, $\dfrac{1}{\sinh^{2}x}$ on $(0,\infty)$. Any denominator that is $1$ minus something small.
\method Use $\dfrac{1}{1-u}=\sum_{k\geq0}u^{k}$ for $\abs{u}<1$. On the half line the useful forms are
\[\frac{1}{\eu^{cx}-1}=\sum_{n\geq1}\eu^{-cnx},\qquad \frac{1}{\eu^{cx}+1}=\sum_{n\geq1}(-1)^{n-1}\eu^{-cnx},\qquad \frac{1}{\sinh^{2}u}=4\sum_{n\geq1}n\,\eu^{-2nu}.\]
Integrate termwise. On $[0,1]$ each term is a power, on $(0,\infty)$ each term is a Gamma moment $\int_0^\infty x^{s-1}\eu^{-nx}\dd x=\Gamma(s)n^{-s}$.

\begin{worked}
Two instances, one on each domain. First,
\[\int_0^1\frac{-\ln(1-x)}{x}\dd x=\int_0^1\sum_{n\geq1}\frac{x^{n-1}}{n}\dd x=\sum_{n\geq1}\frac{1}{n}\cdot\frac1n=\zeta(2)=\frac{\pi^{2}}{6}\approx1.644934.\]
Every term is nonnegative on $[0,1]$, so monotone convergence licenses the swap outright; no epsilon argument is needed.

Second, the Bose integral. Expanding $\frac{1}{\eu^{x}-1}=\sum_{n\geq1}\eu^{-nx}$ and using $\int_0^\infty x^{3}\eu^{-nx}\dd x=\frac{6}{n^{4}}$,
\[\int_0^{\infty}\frac{x^{3}}{\eu^{x}-1}\dd x=6\sum_{n\geq1}\frac{1}{n^{4}}=6\,\zeta(4)=\frac{\pi^{4}}{15}\approx6.493939.\]
The general statement is $\int_0^\infty\frac{x^{s-1}}{\eu^{x}-1}\dd x=\Gamma(s)\zeta(s)$, and with the alternating expansion $\int_0^\infty\frac{x^{s-1}}{\eu^{x}+1}\dd x=\Gamma(s)\eta(s)$; at $s=2$ the latter gives $\eta(2)=\tfrac{\pi^{2}}{12}\approx0.822467$.

The hyperbolic expansions are the same idea after one algebraic step. From $\frac{1}{\sinh u}=\frac{2\eu^{-u}}{1-\eu^{-2u}}=2\sum_{k\geq0}\eu^{-(2k+1)u}$ and $\int_0^\infty u\,\eu^{-(2k+1)u}\dd u=(2k+1)^{-2}$,
\[\int_0^{\infty}\frac{u}{\sinh u}\dd u=2\sum_{k\geq0}\frac{1}{(2k+1)^{2}}=2\cdot\frac34\zeta(2)=\frac{\pi^{2}}{4}\approx2.467401,\]
and differentiating the same expansion in $u$ gives $\frac{1}{\sinh^{2}u}=4\sum_{n\geq1}n\,\eu^{-2nu}$, whence $\int_0^\infty\frac{u^{2}}{\sinh^{2}u}\dd u=4\sum_{n\geq1}n\cdot\frac{1}{4n^{3}}=\zeta(2)$. Both are confirmed numerically.
\end{worked}

\begin{trap}
The expansion fails at the endpoint $\abs{u}=1$, and on $[0,1]$ that endpoint is in the domain. When the terms are all nonnegative this costs nothing, because monotone convergence does not care about a single point. When they alternate it can cost everything, and the correct route is to work on $[0,1-\varepsilon]$, bound the tail, and let $\varepsilon\to0$.
\end{trap}

\begin{keynote}[The three moments]
Termwise integration is only as good as your ability to do the general term, and in practice three moments cover almost everything. On $[0,1]$, $\int_0^1x^{a}\ln^{b}x\dd x=\frac{(-1)^{b}b!}{(a+1)^{b+1}}$. On $(0,\infty)$, $\int_0^\infty x^{s-1}\eu^{-nx}\dd x=\Gamma(s)n^{-s}$. On $[0,\tfrac{\pi}{2}]$, $2\int_0^{\pi/2}\sin^{2p-1}x\cos^{2q-1}x\dd x=\Bfun(p,q)$. Before expanding, decide which of the three your general term will be; if it is none of them, the expansion is probably the wrong move and a substitution should come first.
\end{keynote}

\begin{seealsob}Termwise integration of a Maclaurin series; the log power moment; justifying the swap.\end{seealsob}

Geometric expansion handles denominators. For everything else the expandable factor is a named Maclaurin series, and the difficulty moves from convergence to bookkeeping: matching the index of the series against the index of the resulting sum.

\tech{Termwise integration of a Maclaurin series}{medium}
\cue An integrand with no accessible antiderivative but a standard series: $\arctan x$, $\arcsin x$, $\ln(1+x)$, $\eu^{x\ln x}$, or a Cauchy product such as $(\arctan x)^{2}$.
\method Expand, integrate each power against whatever else is present, and recognise the sum. The self referential case is worth isolating: $x^{x}=\eu^{x\ln x}=\sum_{k\geq0}\frac{(x\ln x)^{k}}{k!}$ turns a transcendental integrand into log power moments.

\begin{worked}
The Sophomore's dream. Expanding and applying $\int_0^1x^{k}\ln^{k}x\dd x=\frac{(-1)^{k}k!}{(k+1)^{k+1}}$,
\[\int_0^1x^{x}\dd x=\sum_{k\geq0}\frac{1}{k!}\cdot\frac{(-1)^{k}k!}{(k+1)^{k+1}}=\sum_{k\geq0}\frac{(-1)^{k}}{(k+1)^{k+1}}=\sum_{n\geq1}\frac{(-1)^{n-1}}{n^{n}}\approx0.783431,\]
and the same computation with $x^{-x}$ flips every sign, giving $\int_0^1x^{-x}\dd x=\sum_{n\geq1}n^{-n}\approx1.291286. $ The reindexing $n=k+1$ is the only place to go wrong, and it is where most people do.

A second instance, with Catalan's constant as the payoff:
\[\int_0^1\frac{\arctan x}{x}\dd x=\sum_{k\geq0}\frac{(-1)^{k}}{2k+1}\int_0^1x^{2k}\dd x=\sum_{k\geq0}\frac{(-1)^{k}}{(2k+1)^{2}}=G\approx0.915966.\]
Here the terms alternate, so monotone convergence is unavailable; the licence comes from the alternating series bound $\bigl\lvert\sum_{k\leq K}\frac{(-1)^{k}x^{2k}}{2k+1}\bigr\rvert\leq1$ on $[0,1]$ uniformly in $K$, which is an integrable dominating function.

A binomial expansion works the same way, with $\binom{2n}{n}$ in place of a factorial. From $(1-u)^{-1/2}=\sum_{n\geq0}\binom{2n}{n}\frac{u^{n}}{4^{n}}$ one gets $\arcsin x=\sum_{n\geq0}\frac{\binom{2n}{n}}{4^{n}(2n+1)}x^{2n+1}$, so
\[\int_0^1\frac{\arcsin x}{x}\dd x=\sum_{n\geq0}\frac{\binom{2n}{n}}{4^{n}(2n+1)^{2}}=\frac{\pi}{2}\ln2\approx1.088793.\]
The closed form is easier to obtain by parts, since $\int_0^1\frac{\arcsin x}{x}\dd x=-\int_0^1\frac{\ln x}{\sqrt{1-x^{2}}}\dd x$ and the right side is a Beta derivative, but the series is the check: summing forty terms reproduces $1.08879$. That pairing is typical, and it is the practical reason to know both routes.
\end{worked}

\begin{keynote}[The same dream on the other domain]
The Sophomore's dream can also be done with the Gamma moment, and comparing the two routes is a good test of whether the machinery is understood. Substituting $x=\eu^{-u}$ carries $[0,1]$ to $(0,\infty)$ and turns $x^{x}\dd x$ into $\eu^{-u\eu^{-u}}\eu^{-u}\dd u$. Expanding the outer exponential and using $\int_0^\infty u^{k}\eu^{-(k+1)u}\dd u=\frac{k!}{(k+1)^{k+1}}$,
\[\int_0^1x^{x}\dd x=\sum_{k\geq0}\frac{(-1)^{k}}{k!}\cdot\frac{k!}{(k+1)^{k+1}}=\sum_{k\geq0}\frac{(-1)^{k}}{(k+1)^{k+1}},\]
the same series with the same index shift. The factorials cancel in both routes, one against the log power moment and one against the Gamma moment, which is not a coincidence: the substitution $x=\eu^{-u}$ is precisely what carries one moment into the other. The series itself converges ferociously, since $n^{-n}$ beats any geometric rate: six terms already give $0.7834305$, correct to seven places.
\end{keynote}

\begin{trap}
Mis-indexing when the $k$th series term produces the $(k+1)$th sum term, as in the Sophomore's dream. Write both indices down explicitly rather than trusting the pattern. The second trap is expanding a factor whose series only converges on part of the interval of integration and then integrating over all of it.
\end{trap}

\begin{seealsob}Geometric expansion; the log power moment; recognising the emergent constant.\end{seealsob}

\subsection{Expansions on the half line}

On $(0,\infty)$ the expandable object is almost always a function of $\eu^{-x}$, and the moment is the Gamma moment. Two immediate consequences, each one line:
\[\int_0^{\infty}\ln\bigl(1-\eu^{-x}\bigr)\dd x=-\sum_{n\geq1}\frac1n\int_0^\infty \eu^{-nx}\dd x=-\zeta(2)=-\frac{\pi^{2}}{6},\]
\[\int_0^{\infty}\ln\bigl(1+\eu^{-x}\bigr)\dd x=\sum_{n\geq1}\frac{(-1)^{n-1}}{n^{2}}=\eta(2)=\frac{\pi^{2}}{12}\approx0.822467.\]
The first has terms of one sign and needs nothing; the second alternates and needs the uniform bound on partial sums, which the alternating series estimate supplies. Raising the power of $x$ raises the weight in the expected way: $\int_0^\infty\frac{x^{2}}{\eu^{x}+1}\dd x=\Gamma(3)\eta(3)=2\cdot\tfrac34\zeta(3)=\tfrac32\zeta(3)\approx1.803085$.

A useful habit on this domain is to substitute $u=\eu^{-x}$ before doing anything else. It carries $(0,\infty)$ to $(0,1)$, turns $\dd x$ into $-\tfrac{\dd u}{u}$, and converts a Bose factor into $\frac{1}{1-u}$, so a half line problem becomes a unit interval problem and the Gamma moment becomes the log power moment. Whichever domain you find easier, one substitution takes you there.

\begin{keynote}[When not to expand]
Expansion is not free, and three signs say to try something else first. If the integrand has a symmetry, use it: $\int_0^{\pi/2}\frac{\dd x}{1+\tan^{\sqrt2}x}$ falls to the reflection $x\mapsto\tfrac{\pi}{2}-x$ in one line and to no series at all. If a substitution makes the integrand rational, make it: partial fractions give an exact answer where a series gives an approximation you then have to recognise. And if the expansion produces a general term you cannot integrate, stop; a series whose terms are as hard as the original problem has cost you a convergence obligation and bought you nothing.
\end{keynote}

\section{The moment that does the work}

Every termwise expansion on $[0,1]$ eventually asks for the integral of a power times a power of a logarithm. That one formula is the engine of the whole chapter, and it is worth knowing not as a lookup but as a derivation, because the derivation tells you the sign.

\tech{The log power moment}{easy}
\cue $x^{a}\ln^{b}x$ anywhere on $[0,1]$, typically appearing after a series expansion has been performed and the general term written out.
\method
\[\int_0^1x^{a}\ln^{b}x\dd x=\frac{(-1)^{b}\,b!}{(a+1)^{b+1}},\qquad a>-1,\ b\in\N_0.\]
The fastest derivation is parameter differentiation: $\int_0^1x^{a}\dd x=\frac{1}{a+1}$, and $\partial_a^{b}$ inserts $\ln^{b}x$ on the left and produces $\frac{(-1)^{b}b!}{(a+1)^{b+1}}$ on the right. The sign is therefore not something to memorise; it is $b$ differentiations of a reciprocal.

\begin{worked}
Directly: $\int_0^1x^{2}\ln^{3}x\dd x=\frac{(-1)^{3}3!}{3^{4}}=-\frac{6}{81}=-\frac{2}{27}\approx-0.074074$.

In use, with $b=2$ and $a=2n$, the moment is $\frac{2}{(2n+1)^{3}}$, which is the engine behind every $\zeta(3)$ that appears in this subject:
\[\int_0^1\frac{\ln^{2}x}{1-x^{2}}\dd x=\sum_{n\geq0}\int_0^1x^{2n}\ln^{2}x\dd x=2\sum_{n\geq0}\frac{1}{(2n+1)^{3}}=2\cdot\frac78\zeta(3)=\frac{7}{4}\zeta(3)\approx2.103600.\]
The factor $\tfrac78$ is the odd index part of $\zeta(3)$, namely $(1-2^{-3})\zeta(3)$. All terms are nonnegative, so the swap is licensed by monotone convergence.

The same moment against the three standard denominators gives a small table worth carrying:
\[\int_0^1\frac{\ln^{2}x}{1-x}\dd x=2\zeta(3),\qquad \int_0^1\frac{\ln^{2}x}{1+x}\dd x=2\eta(3)=\frac32\zeta(3),\qquad \int_0^1\frac{\ln^{2}x}{1+x^{2}}\dd x=2\beta(3)=\frac{\pi^{3}}{16},\]
with numerical values $2.404114$, $1.803085$ and $1.937892$. The three differ only in which subfamily of integers the sum runs over, and reading off the right one is the entire difficulty.
\end{worked}

\begin{trap}
Applying the formula with $a\leq-1$, where the integral diverges: $\int_0^1x^{-1}\ln x\dd x$ is not $-1$, it is infinite. And on $[0,1]$ the logarithm is negative, so an odd $b$ gives a negative answer; a positive result for odd $b$ is a sign error, not a surprise.
\end{trap}

\begin{keynote}[The weight ladder]
Each extra logarithm in the integrand raises the weight of the answer by one. Weight $1$ answers are $\ln2$ and $\tfrac{\pi}{4}$; weight $2$ answers are $\zeta(2)$, $\eta(2)$, $G$ and $\ln^{2}2$; weight $3$ answers are $\zeta(3)$, $\beta(3)$, $\pi^{2}\ln2$ and $\ln^{3}2$; weight $4$ answers are $\zeta(4)$, $\Li_4(\tfrac12)$ and their companions. The rule is not a theorem, but it is a reliable audit: an integral with two logarithms whose reported answer is a weight $2$ constant has lost a factor somewhere, and one whose answer mixes weights is almost certainly wrong, since the pieces of a genuine evaluation are homogeneous.
\end{keynote}

\begin{seealsob}Termwise integration of a Maclaurin series; geometric expansion; recognising the emergent constant.\end{seealsob}

\section{Justifying the swap}

Everything above contained a step of the form $\int\sum=\sum\int$. That step is a theorem, not a notation, and it is false in general. This section states the three usable theorems and then does something more useful: it exhibits two concrete failures, one for a series and one for a double integral, in which the unjustified interchange returns a perfectly finite number that is simply wrong.

It helps to know what the theorems are actually saying. All three are statements about a double integral over a product space, one of whose factors may be a counting measure on the integers, in which case ``integral'' means ``sum''. Tonelli says that for a nonnegative integrand the double integral and both iterated integrals are the same extended real number, with no finiteness hypothesis at all. Fubini says that if the double integral of the absolute value is finite, the same conclusion holds for the signed integrand. Dominated convergence is the version for a limit of partial sums rather than a sum, and it is what you use when the natural object is $\lim_N\int\sum_{n\leq N}$. In every case the hypothesis is checked on $\abs{f}$, never on $f$, and that is the whole content.

\tech{Justifying the swap}{core}
\cue Any step that moves an integral sign past a summation sign, or that reverses the order of two integrals. If you wrote one, you owe a reason.
\method Three tools, in increasing cost. \emph{Monotone convergence}: if $f_n\geq0$, then $\int\sum f_n=\sum\int f_n$ unconditionally, both sides possibly infinite. \emph{Tonelli}: if the double integrand is nonnegative and measurable, both iterated integrals agree with the double integral, again unconditionally. \emph{Fubini and dominated convergence}, for signed integrands: you must first show $\sum\int\abs{f_n}<\infty$, or exhibit an integrable $g$ with $\bigl\lvert\sum_{n\leq N}f_n\bigr\rvert\leq g$ for every $N$. There is no fourth tool, and there is no version in which the swap is automatic.

\begin{worked}
A series where the swap fails and gives a finite wrong answer. On $[0,1]$ put $g_n(x)=n^{2}x\,\eu^{-n^{2}x^{2}}$ and $f_n=g_n-g_{n+1}$. For each fixed $x>0$ the partial sums telescope, $\sum_{n=1}^{N}f_n(x)=x\eu^{-x^{2}}-(N+1)^{2}x\eu^{-(N+1)^{2}x^{2}}\to x\eu^{-x^{2}}$, so
\[\int_0^1\sum_{n\geq1}f_n(x)\dd x=\int_0^1x\eu^{-x^{2}}\dd x=\frac{1-\eu^{-1}}{2}\approx0.316060.\]
Now integrate first. Since $\int_0^1 g_n=\tfrac12\bigl(1-\eu^{-n^{2}}\bigr)$, we get $\int_0^1f_n=\tfrac12\bigl(\eu^{-(n+1)^{2}}-\eu^{-n^{2}}\bigr)$, and that telescopes to
\[\sum_{n\geq1}\int_0^1f_n(x)\dd x=-\frac{\eu^{-1}}{2}\approx-0.183940.\]
The two answers differ, and not by a little: they differ in sign. Numerically both values are exactly reproduced by quadrature. What went wrong is visible: $g_n$ is a spike of fixed mass $\tfrac12$ that migrates to $x=0$, so no integrable $g$ dominates the partial sums and dominated convergence does not apply.

The same phenomenon for a double integral. With $f(x,y)=\dfrac{x^{2}-y^{2}}{(x^{2}+y^{2})^{2}}=\partial_y\dfrac{y}{x^{2}+y^{2}}$ on the unit square,
\[\int_0^1\!\!\left(\int_0^1 f\dd y\right)\dd x=\int_0^1\frac{\dd x}{1+x^{2}}=\frac{\pi}{4},\qquad \int_0^1\!\!\left(\int_0^1 f\dd x\right)\dd y=-\int_0^1\frac{\dd y}{1+y^{2}}=-\frac{\pi}{4},\]
the second following from the first by the antisymmetry $f(y,x)=-f(x,y)$. Both iterated integrals exist and are finite; they disagree because $\iint\abs{f}=\infty$, so Fubini never applied.

The discrete version is even shorter, and worth carrying because it is impossible to misremember. Let $a_{nn}=1$, $a_{n,n+1}=-1$, and $a_{mn}=0$ otherwise. Summing along rows, every row sums to $1-1=0$, so $\sum_m\sum_n a_{mn}=0$. Summing along columns, column $1$ contains only $a_{11}=1$ and every later column contains $a_{nn}=1$ and $a_{n-1,n}=-1$, so $\sum_n\sum_m a_{mn}=1$. Two orders, two answers, and the array is as simple as an array can be.
\end{worked}

\begin{keynote}[What to actually check]
In practice you run one of three checks and record which. If the terms have one sign on the domain, write ``all terms nonnegative, monotone convergence'' and move on. If they alternate, look for a uniform bound on the partial sums, which for an alternating series with decreasing terms is just the first term. If neither holds, compute $\sum\int\abs{f_n}$; if it diverges, do not swap, and instead work on a subinterval where the tail can be estimated by hand.
\end{keynote}

\begin{keynote}[Abel's theorem is not an interchange theorem]
Abel's theorem says that if $\sum a_n$ converges then $\sum a_nx^{n}\to\sum a_n$ as $x\to1^{-}$. It is constantly invoked in this subject and it is constantly misused, because it says nothing whatever about exchanging an integral with a sum. What it does is repair the endpoint after a correct interchange on $[0,1-\varepsilon]$: you prove the identity on the open interval, where uniform convergence on compact subsets is available, and then let $x\to1^{-}$ on both sides, using Abel on the series side and dominated or monotone convergence on the integral side. Two theorems, two jobs. Citing Abel in place of the interchange itself is the most common piece of hand waving in the literature on these integrals.
\end{keynote}

\begin{trap}
Treating a convergent answer as evidence that the swap was legal. Both failures above return finite numbers. The only evidence is a theorem, or a numerical check of the original integral against the summed answer, which is why step 5 of the protocol ends with a numerical comparison.
\end{trap}

\begin{seealsob}Geometric expansion; double integral and double series; convergence tests for series.\end{seealsob}

\begin{keynote}[The finite sum with an explicit remainder]
When the terms alternate and the endpoint is in the domain, there is a technique that never needs a convergence theorem at all. Write the geometric series as a finite sum plus an exact remainder,
\[\frac{1}{1+x}=\sum_{k=0}^{N-1}(-x)^{k}+\frac{(-x)^{N}}{1+x},\]
integrate the identity, which is legal because it is a statement about finitely many terms, and bound the remainder: $\bigl\lvert\int_0^1\frac{(-x)^{N}}{1+x}\dd x\bigr\rvert\leq\int_0^1x^{N}\dd x=\frac{1}{N+1}\to0$. That gives $\sum_{k\geq0}\frac{(-1)^{k}}{k+1}=\ln2$ with no appeal to dominated convergence, Abel's theorem, or uniform convergence. The same device works for every geometric expansion on $[0,1]$ and for most Maclaurin series with an explicit Lagrange remainder. It is more work than quoting a theorem and it is completely airtight, which makes it the right move whenever the theorem is in doubt.
\end{keynote}

\section{Two integrals are easier than one}

The interchange theorems are usually thought of as a tax on an expansion. They are better thought of as a tool in their own right: if swapping the order of two integrals is what makes a problem easy, then a problem with only one integral in it should be given a second one on purpose. The representations of Chapter 23 supply the second integral; this section supplies the reason to want it.

There is a second reason to want the extra dimension, independent of any swap. A double integral over a symmetric region can be averaged with its own reflection, and the average is often elementary where neither half is. If $R$ is symmetric under $(x,y)\mapsto(y,x)$ then $\iint_R f=\tfrac12\iint_R\bigl[f(x,y)+f(y,x)\bigr]$, and the symmetrised integrand frequently collapses: $\frac{x}{x+y}$ becomes $\tfrac12$, and $\frac{\ln x}{\ln x+\ln y}$ becomes $\tfrac12$ likewise. Lifting a one dimensional problem into two dimensions in order to symmetrise it is a distinct move from lifting it in order to swap, and the two are worth keeping separate in your head even though both begin the same way.

\tech{Double integral and double series}{hard}
\cue A single integral that resists everything, but whose integrand is a $y$ integral in disguise: $\ln(1-x)$, $\tfrac1x$, $\tfrac{x^{a}-x^{b}}{\ln x}$, $\arctan x$. Also any integral over $[0,1]^{2}$ of $\tfrac{1}{1-xy}$ shape.
\method Lift to two dimensions, then do one of three things. Swap the order and perform the easy integral first. Expand geometrically into a double series and sum. Or change variables to decouple the region, for which the workhorse maps are $u=x+y$, $v=\tfrac{y}{x+y}$ with Jacobian $u$, and Calabi's map $x=\tfrac{\sin u}{\cos v}$, $y=\tfrac{\sin v}{\cos u}$ sending an open triangle onto the unit square.

\begin{worked}
Start with the deliberate lift. To evaluate $\int_0^1\frac{x-1}{\ln x}\dd x$, note that $\frac{x-1}{\ln x}=\int_0^1x^{t}\dd t$, since the $t$ integral of $x^{t}=\eu^{t\ln x}$ is $\frac{x-1}{\ln x}$. The integrand is nonnegative on $[0,1]^{2}$, so Tonelli permits the swap:
\[\int_0^1\!\!\int_0^1x^{t}\dd t\dd x=\int_0^1\!\!\int_0^1x^{t}\dd x\dd t=\int_0^1\frac{\dd t}{t+1}=\ln2.\]
The same integral was done in Chapter 23 by differentiating under the sign; here nothing was differentiated at all.

Now the classical square. Expanding geometrically,
\[\int_0^1\!\!\int_0^1\frac{\dd x\dd y}{1-xy}=\sum_{n\geq0}\int_0^1\!\!\int_0^1(xy)^{n}\dd x\dd y=\sum_{n\geq0}\frac{1}{(n+1)^{2}}=\zeta(2)=\frac{\pi^{2}}{6},\]
legitimate because every term is positive. Restricting to odd powers gives $\int_0^1\!\!\int_0^1\frac{\dd x\dd y}{1-x^{2}y^{2}}=\tfrac34\zeta(2)=\tfrac{\pi^{2}}{8}\approx1.233701$, and Calabi's map turns that double integral into the area of the triangle $u,v>0$, $u+v<\tfrac{\pi}{2}$, which is $\tfrac{\pi^{2}}{8}$ directly. One level up, $\frac{-\ln(1-u)}{1-u}=\sum_{n\geq1}H_nu^{n}$ with $u=xy$ gives
\[\int_0^1\!\!\int_0^1\frac{-\ln(1-xy)}{1-xy}\dd x\dd y=\sum_{n\geq1}\frac{H_n}{(n+1)^{2}}=\sum_{m\geq1}\frac{H_m}{m^{2}}-\zeta(3)=2\zeta(3)-\zeta(3)=\zeta(3)\approx1.202057,\]
using the Euler sum $\sum H_m m^{-2}=2\zeta(3)$ and the shift $H_{m-1}=H_m-\tfrac1m$. Note that the answer is $\zeta(3)$ and not $2\zeta(3)$: the shift of the index by one is exactly what removes the extra $\zeta(3)$, and dropping it is the standard way to be wrong by a factor of two here.

The lift also proves Frullani. With $\frac1x=\int_0^\infty \eu^{-xt}\dd t$ and $a,b>0$,
\[\int_0^{\infty}\frac{\eu^{-ax}-\eu^{-bx}}{x}\dd x=\int_0^{\infty}\!\!\int_0^{\infty}\Bigl(\eu^{-(a+t)x}-\eu^{-(b+t)x}\Bigr)\dd t\dd x=\int_0^{\infty}\left(\frac{1}{a+t}-\frac{1}{b+t}\right)\dd t=\ln\frac{b}{a},\]
the last integral being a convergent difference of two divergent logarithms, evaluated as $\lim_{T\to\infty}\ln\frac{(a+T)b}{(b+T)a}$. At $a=1$, $b=3$ this is $\ln3\approx1.098612$. The double integrand is not of one sign, but it is absolutely integrable, so Fubini applies; that check takes one line and is the difference between a proof and an assertion.
\end{worked}

\begin{trap}
Getting the transformed \emph{region} wrong. The Jacobian is the easy half of a change of variables; drawing the image of the boundary is where the technique actually fails, and a wrong region gives a wrong number with no symptom. The second trap is lifting an integral that is only conditionally convergent, such as $\int_0^\infty\frac{\sin x}{x}\dd x$ via $\frac1x=\int_0^\infty \eu^{-xt}\dd t$: the swap gives the correct $\int_0^\infty\frac{\dd t}{1+t^{2}}=\frac{\pi}{2}$, but $\iint\abs{\sin x}\eu^{-xt}$ is infinite, so Fubini does not license it and a separate limiting argument is required.
\end{trap}

\begin{keynote}[Symmetrise, then integrate]
The averaging move is worth stating as a formula because it is so easy to forget it exists. On the unit square,
\[\int_0^1\!\!\int_0^1 f(x,y)\dd x\dd y=\frac12\int_0^1\!\!\int_0^1\bigl[f(x,y)+f(y,x)\bigr]\dd x\dd y,\]
which is an identity, not a theorem, since the two orders integrate the same function over the same region. Its value is that the bracket can be far simpler than either term: with $f=\frac{x^{a}}{x^{a}+y^{a}}$ the bracket is identically $1$ and the answer is $\tfrac12$ for every $a$, a result that is invisible from either half alone.
\end{keynote}

\begin{seealsob}Schwinger parametrisation; justifying the swap; symmetrisation by swapping variables.\end{seealsob}

\begin{keynote}[The second integral, from stock]
Manufacturing the extra integral is a lookup, not an invention. The four that recur are
\[\frac1n=\int_0^1t^{n-1}\dd t,\qquad \frac{1}{n^{2}}=-\int_0^1t^{n-1}\ln t\dd t,\qquad \frac{1}{n^{s}}=\frac{1}{\Gamma(s)}\int_0^{\infty}t^{s-1}\eu^{-nt}\dd t,\qquad H_n=\int_0^1\frac{1-t^{n}}{1-t}\dd t.\]
Substituting one of these into a sum and swapping turns the sum geometric, which is the whole trick behind evaluating $\sum\frac{(-1)^{n}}{3n+1}$ and its relatives. Running the same table backwards turns an integral into a sum. Which direction you use depends only on which of the two objects you would rather be holding at the end.
\end{keynote}

\subsection{A lift that needs a limit}

The honest treatment of one famous case is worth writing out, because it is the standard example of a swap that is correct but not licensed by Fubini. Using $\frac1x=\int_0^\infty \eu^{-xt}\dd t$ on $\int_0^\infty\frac{\sin x}{x}\dd x$ produces $\int_0^\infty\frac{\dd t}{1+t^{2}}=\frac{\pi}{2}$, and the answer is right, but $\iint\abs{\sin x}\eu^{-xt}\dd t\dd x=\int_0^\infty\frac{\abs{\sin x}}{x}\dd x=\infty$, so the theorem does not apply.

The repair is to swap on a finite $x$ range, where absolute integrability does hold, and take the limit afterwards. For fixed $R$,
\[\int_0^{R}\frac{\sin x}{x}\dd x=\int_0^{\infty}\left[\int_0^{R}\eu^{-xt}\sin x\dd x\right]\dd t=\int_0^{\infty}\frac{1-\eu^{-Rt}\bigl(\cos R+t\sin R\bigr)}{1+t^{2}}\dd t,\]
the swap now being legitimate because $\int_0^R\int_0^\infty\abs{\sin x}\eu^{-xt}\dd t\dd x=\int_0^R\frac{\abs{\sin x}}{x}\dd x<\infty$. The correction term is bounded by $\int_0^\infty \eu^{-Rt}(1+t)\dd t=\frac1R+\frac{1}{R^{2}}$, since $\frac{1+t}{1+t^{2}}\leq1+t$, so it vanishes as $R\to\infty$ and
\[\int_0^{\infty}\frac{\sin x}{x}\dd x=\int_0^{\infty}\frac{\dd t}{1+t^{2}}=\frac{\pi}{2}.\]
Three extra lines converted an illegal manipulation into a proof. That is the usual price, and it is almost always this cheap: truncate, swap where the theorem applies, estimate the tail.

\subsection{Running the lift backwards: a binomial sum}

The table above is usually read from the integral towards the sum. Read the other way, it evaluates sums that no summation technique touches. Take
\[\Sigma=\sum_{n\geq1}\frac{1}{n^{2}\binom{2n}{n}}.\]
The Beta function supplies a representation of the reciprocal binomial: since $\Bfun(n,n+1)=\frac{(n-1)!\,n!}{(2n)!}=\frac{1}{n\binom{2n}{n}}$, we have $\frac{1}{\binom{2n}{n}}=n\int_0^1x^{n-1}(1-x)^{n}\dd x$. Substituting and swapping, which is licensed because every term is positive on $[0,1]$,
\[\Sigma=\sum_{n\geq1}\frac1n\int_0^1x^{n-1}(1-x)^{n}\dd x=\int_0^1\frac1x\sum_{n\geq1}\frac{\bigl[x(1-x)\bigr]^{n}}{n}\dd x=-\int_0^1\frac{\ln\bigl(1-x+x^{2}\bigr)}{x}\dd x.\]
The remaining integral is a standard one, and both sides evaluate to $\frac{\pi^{2}}{18}\approx0.548311$, which matches the numerical sum to fifteen places. The companion sum $\sum_{n\geq1}\frac{1}{n\binom{2n}{n}}=\frac{\pi\sqrt3}{9}\approx0.604600$ comes from the same representation with one fewer factor of $\tfrac1n$.

What made this work was that the binomial coefficient, which is an obstruction to summation, is a Beta integral, which is not an obstruction to anything. That is the same observation that drove the whole of the previous chapter, applied to a sum instead of an integral.

\subsection{Changing the region instead of the integrand}

The third route in the method above deserves its own treatment, because it is where a double integral stops being a bookkeeping device and starts being a geometric argument. Restricting the double geometric series to odd powers gives
\[\int_0^1\!\!\int_0^1\frac{\dd x\dd y}{1-x^{2}y^{2}}=\sum_{k\geq0}\frac{1}{(2k+1)^{2}}=\frac34\zeta(2),\]
so evaluating the left side by any means evaluates $\zeta(2)$. Calabi's substitution
\[x=\frac{\sin u}{\cos v},\qquad y=\frac{\sin v}{\cos u}\]
maps the open triangle $u>0$, $v>0$, $u+v<\tfrac{\pi}{2}$ bijectively onto the open unit square, and its Jacobian is exactly $1-x^{2}y^{2}$. The integrand therefore cancels the Jacobian and the double integral is the plain area of the triangle,
\[\frac34\zeta(2)=\frac12\left(\frac{\pi}{2}\right)^{2}=\frac{\pi^{2}}{8}\approx1.233701,\qquad\text{hence}\qquad \zeta(2)=\frac{\pi^{2}}{6}.\]
This is a complete and elementary proof of Euler's identity, and the only analytic step in it is the geometric expansion, which is licensed by positivity. What made it work is that the change of variables was chosen to kill the integrand rather than to simplify the region, which is the opposite of the usual instinct and worth remembering as an option.

The companion map $u=x+y$, $v=\tfrac{y}{x+y}$, with Jacobian $u$, does the reverse: it takes a triangle to a rectangle and separates a symmetric integrand into a product. Between them the two maps handle most integrals over the unit square whose integrand depends on $xy$ or on $x+y$.

\section{Full periods and orthogonality}

On a bounded interval the natural expansion is in powers. On a full period the natural expansion is in cosines, and it has a property powers do not: distinct modes are orthogonal, so a double sum collapses to a single one on integration. This is what makes logarithmic trigonometric integrals tractable.

The two expansions used below are worth deriving once rather than memorising. Both come from the same source: for $\abs{r}<1$,
\[-\ln(1-r\eu^{\iu x})=\sum_{n\geq1}\frac{r^{n}\eu^{\iu nx}}{n},\]
whose real part at $r=1$ is $-\ln\bigl(2\sin\tfrac x2\bigr)$ after the identity $\abs{1-\eu^{\iu x}}=2\sin\tfrac x2$ on $(0,2\pi)$, and whose companion, obtained by summing the geometric series $\sum_{n\geq1}r^{n}\eu^{\iu nx}$ and taking twice the real part plus one, is the Poisson kernel. So one generating function produces both the logarithm and the kernel, which is why they pair so well under an integral.

\tech{Fourier kernels and orthogonality}{hard}
\cue $\ln\bigl(2\sin\tfrac x2\bigr)$ or $\ln\bigl(2\cos\tfrac x2\bigr)$, or a Poisson kernel shape $\dfrac{1}{a-\cos x}$, integrated over a full period.
\method Expand both factors in cosine series and integrate over the period, where $\int_0^{2\pi}\cos(nx)\cos(mx)\dd x=\pi\delta_{nm}$ for $n,m\geq1$ kills every cross term. The two expansions to know are
\[\ln\Bigl(2\sin\frac x2\Bigr)=-\sum_{n\geq1}\frac{\cos nx}{n}\quad\text{on }(0,2\pi),\qquad \frac{1-r^{2}}{1-2r\cos x+r^{2}}=1+2\sum_{n\geq1}r^{n}\cos nx.\]

\begin{worked}
Take $r=\tfrac12$, so $1-2r\cos x+r^{2}=\tfrac54-\cos x$ and $1-r^{2}=\tfrac34$, giving $\dfrac{1}{\tfrac54-\cos x}=\dfrac43\Bigl(1+2\sum_{n\geq1}2^{-n}\cos nx\Bigr)$. Multiply by the logarithm series and integrate over $[0,2\pi]$. The constant $1$ pairs with every $\cos mx$ to give zero, and orthogonality keeps only $n=m$:
\[\int_0^{2\pi}\frac{\ln\bigl(2\sin\tfrac x2\bigr)}{\tfrac54-\cos x}\dd x=\frac43\cdot2\pi\sum_{n\geq1}2^{-n}\left(-\frac1n\right)=-\frac{8\pi}{3}\ln2\approx-5.806896,\]
since $\sum_{n\geq1}\frac{2^{-n}}{n}=\ln2$. Quadrature agrees to thirteen places.

The same expansion against a different weight produces an odd zeta value. Using $\int_0^{\pi}x\cos nx\dd x=\frac{(-1)^{n}-1}{n^{2}}$,
\[\int_0^{\pi}x\ln\Bigl(2\sin\frac x2\Bigr)\dd x=-\sum_{n\geq1}\frac{(-1)^{n}-1}{n^{3}}=2\sum_{n\ \mathrm{odd}}\frac{1}{n^{3}}=\frac{7}{4}\zeta(3)\approx2.103600.\]
Note that this integral runs over half a period, and the cross terms did not need to vanish because the weight was integrated against each mode separately.

The cheapest consequence of the same series is the log sine integral itself. On $(0,\pi)$ the doubled form reads $\ln(2\sin x)=-\sum_{n\geq1}\frac{\cos2nx}{n}$, and $\int_0^{\pi/2}\cos2nx\dd x=\frac{\sin n\pi}{2n}=0$ for every $n\geq1$, so every term integrates to zero and
\[\int_0^{\pi/2}\ln(2\sin x)\dd x=0,\qquad\text{that is}\qquad \int_0^{\pi/2}\ln\sin x\dd x=-\frac{\pi}{2}\ln2\approx-1.088793.\]
The same computation with $\cos$ in place of $\sin$ gives the same value, as the reflection $x\mapsto\tfrac{\pi}{2}-x$ requires. It is worth noticing how little was used: not orthogonality, but the accident that each mode has zero mean over the quarter period.
\end{worked}

Squaring the logarithm turns the single sum into a double sum, and orthogonality is then doing real work rather than bookkeeping. On $(0,\pi)$, $\ln(2\sin x)=-\sum_{n\geq1}\frac{\cos2nx}{n}$, so
\[\int_0^{\pi/2}\ln^{2}(2\sin x)\dd x=\sum_{m,n\geq1}\frac{1}{mn}\int_0^{\pi/2}\cos2mx\cos2nx\dd x=\frac{\pi}{4}\sum_{n\geq1}\frac{1}{n^{2}}=\frac{\pi^{3}}{24}\approx1.291928,\]
because $\int_0^{\pi/2}\cos2mx\cos2nx\dd x=\tfrac{\pi}{4}\delta_{mn}$ on the quarter period: the cross terms integrate to zero there for the same reason the single terms did. The companion over a full half period, $\int_0^{\pi}\ln^{2}\bigl(2\cos\tfrac x2\bigr)\dd x=\tfrac{\pi^{3}}{12}\approx2.583856$, is the same computation rescaled. Notice the weight accounting: two logarithms, and a weight three answer.

\begin{keynote}[When the period is partial]
Over less than a full period the modes are no longer orthogonal in general, and what survives is a named function rather than a constant. The relevant one is the Clausen function $\mathrm{Cl}_2(\theta)=\sum_{n\geq1}\frac{\sin n\theta}{n^{2}}$, which satisfies $\int_0^{\theta}\ln\bigl(2\sin\tfrac x2\bigr)\dd x=-\mathrm{Cl}_2(\theta)$ and takes the values $\mathrm{Cl}_2(\pi)=0$ and $\mathrm{Cl}_2(\tfrac{\pi}{2})=G$. An integral over a partial period that refuses to close is usually a Clausen value in disguise, and reporting it as such is a complete answer, not a failure.
\end{keynote}

\begin{trap}
The $\ln\sin$ series converges only conditionally and only on $(0,2\pi)$. Over a partial period the orthogonality relation fails and cross terms survive, which is where Clausen functions come from. Integrating $\ln\bigl(2\sin\tfrac x2\bigr)$ over $(0,\theta)$ gives $-\mathrm{Cl}_2(\theta)$, not a rational multiple of anything.
\end{trap}

\begin{seealsob}Log sine integrals and the Clausen function; product to sum identities; geometric expansion.\end{seealsob}

\section{Reading the answer}

A termwise integration ends with a series, and the last step is naming it. This is a dictionary problem, and the dictionary is short, but two entries in it are confused so regularly that the confusion deserves its own warning.

Naming is not merely cosmetic. A closed form is checkable, comparable against other derivations, and usable as an input to the next problem, whereas an unnamed series is none of those things. It is also the step at which an error becomes visible: a wrong factor of two in the integration is invisible in a series but obvious when the series is $\tfrac{\pi^{2}}{12}$ and you expected $\tfrac{\pi^{2}}{6}$. Learn the dictionary in both directions, from series to constant and from constant back to the integral that generates it, because recognition failures in this subject are almost always failures in the second direction.

\tech{Recognising the emergent constant}{easy}
\cue A sum of the form $\sum n^{-s}$, $\sum(-1)^{n-1}n^{-s}$, $\sum(-1)^{k}(2k+1)^{-s}$, $\sum(2n+1)^{-s}$, or $\sum\frac{1}{n(n+1)}$ left standing at the end of a computation.
\method Know the dictionary:
\[\zeta(2)=\frac{\pi^{2}}{6},\quad \zeta(4)=\frac{\pi^{4}}{90},\quad \eta(s)=\bigl(1-2^{1-s}\bigr)\zeta(s),\quad \sum_{n\ \mathrm{odd}}n^{-s}=\bigl(1-2^{-s}\bigr)\zeta(s),\]
\[\eta(1)=\ln2,\quad \eta(2)=\frac{\pi^{2}}{12},\quad \beta(2)=G\approx0.915966,\quad \Li_4(1)-\Li_4(-1)=\frac{15}{8}\zeta(4).\]
No closed form is known for $\zeta(3)$ or any odd zeta value, so landing on one is information, not an error.

\begin{worked}
\[\int_0^1\frac{\ln x\,\ln(1-x)}{x}\dd x.\]
Expand the second logarithm: $-\ln(1-x)=\sum_{n\geq1}\frac{x^{n}}{n}$, so the integrand is $-\sum_{n\geq1}\frac{x^{n-1}\ln x}{n}$. The log power moment with $b=1$ gives $\int_0^1x^{n-1}\ln x\dd x=-\frac{1}{n^{2}}$, hence
\[\int_0^1\frac{\ln x\ln(1-x)}{x}\dd x=\sum_{n\geq1}\frac{1}{n}\cdot\frac{1}{n^{2}}=\zeta(3)\approx1.202057.\]
Every term is nonnegative on $[0,1]$, since $\ln x$ and $\ln(1-x)$ are both negative there.

An alternating cousin, to fix the $\zeta$ against $\eta$ distinction: $\frac{1}{1+x}=\sum_{n\geq0}(-1)^{n}x^{n}$ gives
\[\int_0^1\frac{\ln x}{1+x}\dd x=\sum_{n\geq0}(-1)^{n}\left(-\frac{1}{(n+1)^{2}}\right)=-\eta(2)=-\frac{\pi^{2}}{12}\approx-0.822467,\]
which is exactly half of $-\zeta(2)$ and not $-\zeta(2)$ itself. One weight higher, the same expansion with $\ln^{2}x$ gives $\int_0^1\frac{\ln^{2}x}{1+x}\dd x=2\eta(3)=\tfrac32\zeta(3)\approx1.803085$, and the appearance of an odd zeta value there is not a failure: no closed form exists, and $\tfrac32\zeta(3)$ is the answer.
\end{worked}

\begin{trap}
Reporting $\eta(s)$ as $\zeta(s)$. The alternating sum is $\eta(2)=\tfrac{\pi^{2}}{12}$, half of $\zeta(2)$, and the factor of two is the single most common wrong answer in this chapter. The companion slip is writing $\beta(2)=\tfrac{\pi^{2}}{8}$: the alternating sum over odd squares is Catalan's constant $G$, which is not known to be a rational multiple of $\pi^{2}$, whereas the non alternating sum over odd squares is $\tfrac{\pi^{2}}{8}$.
\end{trap}

\begin{keynote}[The polylogarithm as a bookkeeping device]
When the sum does not collapse to a named constant, name it anyway. $\Li_s(z)=\sum_{n\geq1}\frac{z^{n}}{n^{s}}$ specialises to $\Li_s(1)=\zeta(s)$, $\Li_s(-1)=-\eta(s)$, $\Li_1(z)=-\ln(1-z)$, and $\Li_2(\tfrac12)=\tfrac{\pi^{2}}{12}-\tfrac{\ln^{2}2}{2}$. Two identities do most of the work: the Landen reflection $\Li_2(z)+\Li_2(1-z)=\tfrac{\pi^{2}}{6}-\ln z\ln(1-z)$, and the duplication $\Li_s(z)+\Li_s(-z)=2^{1-s}\Li_s(z^{2})$. An answer expressed in polylogarithms at $\tfrac12$, $-1$ and $1$ is a finished answer; the residual urge to reduce it further is what leads to invented closed forms for $\zeta(3)$.
\end{keynote}

\begin{seealsob}Catalan's constant; the polylogarithm ladder; the log power moment.\end{seealsob}

\subsection{Where each constant comes from}

It is worth knowing the standard integral source of each constant, because recognition then runs in both directions and an unexpected constant becomes a diagnostic rather than a puzzle.

$\ln2$ comes from $\sum\frac{(-1)^{n-1}}{n}$, and therefore from any expansion of $\frac{1}{1+x}$ integrated once on $[0,1]$; the archetype is $\int_0^1\frac{\dd x}{1+x}$.

$\tfrac{\pi}{4}$ comes from $\sum\frac{(-1)^{k}}{2k+1}$, so from $\frac{1}{1+x^{2}}$; a $\pi$ with no logarithm attached almost always entered this way.

$\zeta(2)$ comes from one logarithm against $\frac{1}{1-x}$, as in $\int_0^1\frac{-\ln(1-x)}{x}\dd x$, and its alternating partner $\eta(2)=\tfrac{\pi^{2}}{12}$ from $\int_0^1\frac{-\ln x}{1+x}\dd x$.

Catalan's constant $G=\beta(2)$ comes from one logarithm against $\frac{1}{1+x^{2}}$, that is from odd indices with alternating signs. The two cleanest sources are $\int_0^1\frac{\arctan x}{x}\dd x=G$ and $\int_0^1\frac{\ln x}{1+x^{2}}\dd x=-G$, and the trigonometric form $\int_0^{\pi/4}\ln(\tan x)\dd x=-G$ is the same integral after $x=\tan\theta$. All three agree numerically at $0.915966$.

$\zeta(3)$ comes from two logarithms, whether as $\ln^{2}x$ against $\frac{1}{1-x}$ or as $\ln x\ln(1-x)$ against $\frac1x$, and $\beta(3)=\tfrac{\pi^{3}}{32}$ is its odd alternating cousin.

$\Li_2(\tfrac12)=\tfrac{\pi^{2}}{12}-\tfrac{\ln^{2}2}{2}\approx0.582241$ is what appears when the geometric ratio is $\tfrac12$ rather than $1$, which happens after any substitution that halves an interval.

\begin{keynote}[Sums you have not met]
If the series does not match the dictionary, try three things before giving up. Split by parity, since many sums are a $\zeta$ and an $\eta$ added together. Look for a harmonic number, since $\sum\frac{H_n}{n^{2}}=2\zeta(3)$ and $\sum\frac{H_n}{n^{3}}=\tfrac{\pi^{4}}{72}$ close most Euler sums of low weight. And check whether the general term is a Cauchy product, which it is whenever the integrand contained a square of a logarithm or a product of two expandable factors; recognising the product lets you sum the two factors separately.
\end{keynote}

The final technique is the largest generalisation of expand and integrate: it says that if you know the coefficients of an expansion as a function of the index, the Mellin transform is that same function evaluated at a negative argument. It is powerful, it is often the fastest route, and it is the easiest thing in this chapter to abuse.

\tech{Ramanujan's master theorem}{hard}
\cue A Mellin shaped integral $\int_0^{\infty}x^{s-1}\varphi(x)\dd x$ where $\varphi$ has a Maclaurin expansion whose coefficients are a recognisable function of the index.
\method If $\varphi(x)=\sum_{k\geq0}\frac{\phi(k)}{k!}(-x)^{k}$ with $\phi$ analytic and of controlled growth in a right half plane, then
\[\int_0^{\infty}x^{s-1}\varphi(x)\dd x=\Gamma(s)\,\phi(-s).\]
The content is that you continue the coefficient function to a negative argument rather than summing anything.

\begin{worked}
Two applications, one safe and one delicate. Safe: $\frac{1}{1+x}=\sum_{k\geq0}(-x)^{k}$, so $\phi(k)=k!$, and continuing $k!$ to $\Gamma(k+1)$ gives
\[\int_0^{\infty}\frac{x^{s-1}}{1+x}\dd x=\Gamma(s)\Gamma(1-s)=\frac{\pi}{\sin\pi s},\qquad 0<s<1,\]
which at $s=\tfrac13$ predicts $\pi/\sin\tfrac{\pi}{3}\approx3.627599$, confirmed by quadrature to eight places.

Delicate: $\cos x=\sum_{k\geq0}\frac{(-1)^{k}x^{2k}}{(2k)!}$ has $\phi$ supported on even indices, and the correct continuation is $\phi(k)=\cos\frac{\pi k}{2}$, giving
\[\int_0^{\infty}x^{s-1}\cos x\dd x=\Gamma(s)\cos\frac{\pi s}{2},\qquad 0<s<1.\]
At $s=\tfrac12$ this is $\Gamma(\tfrac12)\cos\tfrac{\pi}{4}=\sqrt{\pi}\cdot\tfrac{\sqrt2}{2}=\sqrt{\tfrac{\pi}{2}}\approx1.253314$, the Fresnel value. The integral converges only conditionally, so the number must be read as a limit of $\int_0^{R}$ along $R\to\infty$; averaging the partial integrals over successive half periods reproduces $1.2533141$ numerically.
\end{worked}

The theorem is best understood as a formal statement made rigorous by a contour argument: if you write $\varphi$ as an inverse Mellin transform and slide the contour, the residues at the poles of $\Gamma$ reproduce the Maclaurin coefficients, and the theorem is the statement that the correspondence runs both ways. That picture explains both its power and its fragility. It is exact when $\phi$ extends to a function with controlled growth in a vertical strip, and it is meaningless when the extension is not unique, which is the usual failure: $\phi(k)=\cos\tfrac{\pi k}{2}$ and $\phi(k)=\cos\tfrac{\pi k}{2}+\sin\pi k$ agree at every nonnegative integer and disagree everywhere else, and only the first satisfies the growth hypothesis. Carlson's theorem is what makes the choice canonical.

\begin{trap}
Applying the theorem without checking growth. The hypotheses require $\phi$ analytic in a strip and bounded by $C\eu^{-\delta\abs{\Im z}}$ type estimates; without them the formula returns confident nonsense, and several classical ``paradoxes'' are exactly this. When the integral is only conditionally convergent, the output is a regularised value and should be labelled as one.
\end{trap}

\begin{seealsob}Euler reflection and the Mellin transform; the Gamma function moment; Fresnel integrals.\end{seealsob}

In practice the theorem earns its place as a fast route to a conjecture rather than as a proof. Apply it, get a candidate closed form in seconds, then verify the candidate numerically and, if it matters, re-derive it by expansion and interchange with the licences written out. That workflow uses the theorem for what it is genuinely good at, which is guessing correctly, and leaves the burden of proof with methods whose hypotheses you can actually check on the page.

\begin{center}
\begin{tikzpicture}[node distance=6mm and 8mm]
\node[dtest] (a) {Is there a $\dfrac{1}{1\pm u}$\\ with $\abs{u}<1$ on the domain?};
\node[dleaf, below left=9mm and 4mm of a] (b) {Geometric series,\\ then a log power\\ or Gamma moment};
\node[dtest, below right=9mm and 4mm of a] (c) {Is some factor a\\ standard Maclaurin series?};
\node[dleaf, below left=9mm and 0mm of c] (d) {Expand it; integrate\\ $x^{a}\ln^{b}x$ termwise};
\node[dtest, below right=9mm and 0mm of c] (e) {Can a factor be\\ written as an integral?};
\node[dleaf, below left=9mm and 0mm of e] (f) {Lift to a double integral\\ and swap the order};
\node[dleaf, below right=9mm and 0mm of e] (g) {Mellin shape: try\\ Ramanujan's theorem};
\draw[dedge] (a) -- node[dlab,left]{yes} (b);
\draw[dedge] (a) -- node[dlab,right]{no} (c);
\draw[dedge] (c) -- node[dlab,left]{yes} (d);
\draw[dedge] (c) -- node[dlab,right]{no} (e);
\draw[dedge] (e) -- node[dlab,left]{yes} (f);
\draw[dedge] (e) -- node[dlab,right]{no} (g);
\end{tikzpicture}
\end{center}

The tree routes on the shape of the integrand, and its branches are ordered by cost rather than by cleverness. Note that the left branch is not always the cheapest: a geometric expansion on $[0,1]$ is free of convergence worries only when the terms keep one sign, and an alternating expansion at the endpoint can cost more work than a direct lift to two dimensions.

\section{Assembling the pieces}

One extended example, to show the protocol running with every part in place. Evaluate
\[S=\int_0^1\frac{\ln^{2}(1-x)}{x}\dd x.\]
Neither factor is expandable by itself in a useful way, but the square of a logarithm is a Cauchy product. From $-\ln(1-x)=\sum_{n\geq1}\frac{x^{n}}{n}$, squaring and collecting the coefficient of $x^{n}$ gives
\[\ln^{2}(1-x)=\sum_{n\geq2}\left(\sum_{j=1}^{n-1}\frac{1}{j(n-j)}\right)x^{n}=\sum_{n\geq2}\frac{2H_{n-1}}{n}x^{n},\]
the inner sum collapsing because $\frac{1}{j(n-j)}=\frac1n\bigl(\frac1j+\frac{1}{n-j}\bigr)$. All coefficients are positive and $x>0$, so monotone convergence licenses the swap without further comment. Integrating termwise against $\frac1x$,
\[S=\sum_{n\geq2}\frac{2H_{n-1}}{n}\cdot\frac1n=2\sum_{n\geq1}\frac{H_n-\tfrac1n}{n^{2}}=2\bigl[2\zeta(3)-\zeta(3)\bigr]=2\zeta(3)\approx2.404114,\]
using Euler's sum $\sum H_nn^{-2}=2\zeta(3)$. Quadrature returns $2.4041138$, so the answer stands.

Three checks were available and all three pass. The weight count: two logarithms plus the $\tfrac1x$ predicts weight $3$, and $\zeta(3)$ has weight $3$. The index shift: $H_{n-1}=H_n-\tfrac1n$ was applied explicitly rather than assumed, which is where the factor of two lives. And the numerical comparison, which costs a few seconds. A calculation that survives all three is almost certainly right; one that survives none of them is not an answer, it is a guess.

For contrast, the closely related $\int_0^1\frac{\ln x\ln(1+x)}{x}\dd x$ expands to $-\sum_{n\geq1}\frac{(-1)^{n-1}}{n}\cdot\frac{1}{n^{2}}=-\eta(3)=-\tfrac34\zeta(3)\approx-0.901543$, and $\int_0^1\frac{\ln(1-x)\ln^{2}x}{x}\dd x=-2\zeta(4)=-\tfrac{\pi^{4}}{45}\approx-2.164646$. The first has an alternating expansion and needs the uniform partial sum bound; the second is of one sign and needs nothing. Same family, different licences.

\begin{keynote}[The order in which to work]
Do the algebra before the analysis. Expand, integrate the general term, sum, and name the constant first, treating every interchange as provisional; then go back and discharge the licences. Working in that order costs nothing, because the licence you need is determined by the expansion you actually used, and it saves the common waste of proving a careful convergence result about a series you then abandon. What it must not become is an excuse to skip the second pass. Write the provisional answer, run the numerical comparison, and only then decide whether the interchange needs monotone convergence, a uniform bound, or a truncation argument.
\end{keynote}

\begin{keynote}[Confirm before you commit]
Every closed form in this chapter can be checked in under a minute: evaluate the original integral numerically and compare with the summed answer. The characteristic error is not a wrong method but a factor of two, arising from an index shift, from confusing $\zeta$ with $\eta$, or from an odd or even restriction applied to the wrong half of a sum. All three show up immediately in the second decimal place. If the numerical value is exactly twice or exactly half your answer, do not look for a subtle analytic mistake; look for the index.
\end{keynote}

The chapter has two halves that look different and are not. Expanding an integrand into a series and lifting an integral into a double integral are the same operation performed against two different measures, counting measure in the first case and Lebesgue measure in the second, and the theorem that licenses each is the same theorem. That is why the failures look alike: a spike of fixed mass escaping to the edge of the domain in the series example, and an unsigned singularity at the corner in the double integral example. Both are failures of absolute integrability, and both are invisible in the final number.

\section{Recognition matrix}
\begin{recog}{@{}p{0.30\textwidth}p{0.36\textwidth}p{0.28\textwidth}@{}}
\rowcolor{mist}\mh{If you see} & \mh{Reach for} & \mh{If that fails} \\
\midrule
\endhead
$\dfrac{1}{1\pm x^{n}}$ on $[0,1]$ & Geometric expansion, then $\int_0^1x^{a}\ln^{b}x$ & Substitute $x^{n}=u$, then Beta \\
$\dfrac{1}{\eu^{x}-1}$ or $\dfrac{1}{\eu^{x}+1}$ on $(0,\infty)$ & $\sum \eu^{-nx}$, giving $\Gamma(s)\zeta(s)$ or $\Gamma(s)\eta(s)$ & Contour integration \\
$\ln(1-x)$, $\arctan x$, $\arcsin x$ over $x$ & Maclaurin series, integrate termwise & Parameter differentiation \\
$x^{x}$, $x^{-x}$, $\eu^{x\ln x}$ & Expand $\eu^{u}$; each term is a log power moment & No closed form; report the series \\
$x^{a}\ln^{b}x$ on $[0,1]$ & $\dfrac{(-1)^{b}b!}{(a+1)^{b+1}}$, with $a>-1$ & Differentiate $\int_0^1x^{a}$ in $a$ \\
A swap of $\int$ and $\sum$ & Sign check first: nonnegative means free & Bound $\sum\int\abs{f_n}$ \\
Two orders of a double integral & Tonelli if nonnegative, else Fubini & Compute both and compare \\
$\dfrac{1}{1-xy}$ on the unit square & Double geometric series, giving $\zeta(2)$ & Change variables to a triangle \\
$\dfrac{x^{a}-x^{b}}{\ln x}$ & Write it as $\int_a^b x^{t}\dd t$ and swap & Differentiate under the sign \\
$\ln\bigl(2\sin\tfrac x2\bigr)$ over a full period & Cosine series and orthogonality & Clausen function on a partial period \\
$\sum(-1)^{n-1}n^{-s}$ & $\eta(s)=(1-2^{1-s})\zeta(s)$ & Check against $\zeta$, factor of two \\
$\sum(-1)^{k}(2k+1)^{-2}$ & Catalan's constant $G$ & Not a multiple of $\pi^{2}$ \\
$\int_0^\infty x^{s-1}\varphi(x)\dd x$, $\varphi$ with nice coefficients & Ramanujan's master theorem, $\Gamma(s)\phi(-s)$ & Verify growth, or expand directly \\
$\ln(1\pm \eu^{-x})$ on $(0,\infty)$ & Expand in $\eu^{-nx}$; Gamma moment & Substitute $u=\eu^{-x}$ \\
$\ln^{2}(1-x)$ or a product of two series & Cauchy product; expect $H_n$ coefficients & Euler sums of weight three \\
Two logarithms in the integrand & Expect weight three: $\zeta(3)$, $\beta(3)$, $\pi^{2}\ln2$ & Recount; a weight mismatch is an error \\
$\dbinom{2n}{n}$ downstairs in a sum & $\Bfun(n,n+1)$ representation, then swap & Generating function of $\binom{2n}{n}$ \\
An alternating expansion at $x=1$ & Finite sum plus explicit remainder & Work on $[0,1-\varepsilon]$, then Abel \\
\end{recog}
